% Part: proof-theory
% Chapter: proof-search
% Section: search-algorithm

\documentclass[../../../include/open-logic-section]{subfiles}

\begin{document}

\olfileid{pt}{ps}{sal}

\olsection{\Log{G3c} కోసం అన్వేషణ అల్గోరిథం}

\Log{G3c} కోసం ఒక అన్వేషణ అల్గోరిథాన్ని వివరిద్దాం. సాధ్యమైనది ఇదొక్కటే కాదు; అత్యంత సమర్థమైనదీ కాదు. అయితే ఇది సరైనది: సీక్వెంట్~$\Gamma \Sequent \Delta$కు \tetoken{ఒక నిరూపణ}{!!a{proof}} ఉంటే, ఇది చివరకు ముగిసి \tetoken{ఒక నిరూపణను}{!!a{proof}} కనుగొంటుంది. ఇది పొరపాటైన నిర్ణయం కూడా ఇవ్వదు: \tetoken{ఒక నిరూపణను}{!!a{proof}} కనుగొనకుండానే ముగిసినా, ఎప్పటికీ ముగియకపోయినా, $\Gamma
\Sequent \Delta$ మొదటి నుంచే సర్వసత్యం కాదు.

అల్గోరిథానికి అవసరమైన కీలక లక్షణం ఇది: $\Gamma \Sequent \Delta$లోని లేదా అన్వేషణలో ఉత్పత్తి అయ్యే సీక్వెంట్‌లోని ప్రతి \tetoken{సూత్రం}{!!{formula}} "తగ్గించబడాలి"; అంటే వెనుకకు వర్తింపజేసిన నిగమన నియమానికి ప్రధాన సూత్రంగా పరిగణించబడాలి; అవసరమైనన్ని సార్లు అలా జరగాలి. ఈ లక్షణాన్ని \emph{సముచిత అవకాశ నిబంధన} అంటాం.

అల్గోరిథం దశలవారీగా పనిచేస్తుంది. ప్రతి దశ ఫలితం సీక్వెంట్‌ల వృక్షం; తరువాతి దశలో, ఇప్పటికే ఆక్సియమ్ కాని ప్రతి అత్యుపరి సీక్వెంట్‌లో \tetoken{ఒక సూత్రాన్ని}{!!a{formula}} తగ్గించి కొత్త ఫలితాన్ని పొందుతాం.

$0$ దశలో $\Gamma \Sequent \Delta$తో మొదలుపెడతాం. సముచిత అవకాశ నిబంధనను నిర్ధారించడానికి, $\Gamma$, $\Delta$లలోని \tetoken{సూత్రాలకు}{!!{formula}s} సంఖ్యలను సూచికలుగా కేటాయిస్తాం; వేర్వేరు \tetoken{సూత్రాలకు}{!!{formula}s} వేర్వేరు సంఖ్యలు వస్తాయి. ఉదాహరణకు \tetoken{సూత్రాలను}{!!{formula}s} ఎడమ నుంచి కుడికి లెక్కపెట్టినట్టు భావించండి. సంఖ్యలను ఊర్ధ్వాంకాలుగా రాస్తాం. $!A, !B \Sequent !C$కు \tetoken{ఒక నిరూపణను}{!!a{proof}} అన్వేషిస్తే,~$0$ దశ ఫలితం $!A^1, !B^2 \Sequent !C^3$. సూచికలతో కూడిన సీక్వెంట్‌లో ఊర్ధ్వాంకంగా కనిపించే అతిపెద్ద సంఖ్య~$n$ను దాని \emph{గరిష్ఠ సూచిక} అంటాం (ఉదాహరణలో~$3$).

సీక్వెంట్‌లో పరిమాణకాలు ఉంటే, కుడివైపున $\lexists[x][!A(x)]$ను లేదా ఎడమవైపున $\lforall[x][!A(x)]$ను తగ్గించేటప్పుడు ఏ పదాలు వాడాలో కూడా తెలియాలి. ~$\Gamma$, $\Delta$లలోని ప్రమేయ చిహ్నాలను ఉపయోగించి \tetoken{స్థిరాంకాలు}{!!{constant}s} $\Obj c_0$, $\Obj c_1$, $\Obj c_2$, \dots నుంచి ఏర్పడే పదాలు మాత్రమే కావాలి. ఈ పదాలన్నింటి సమితికి ఒక స్థిర గణన క్రమం ఉందని అనుకుందాం; అప్పుడు ఉదాహరణకు "$\Gamma \Sequent \Delta$లో కనిపించని మొదటి \tetoken{స్థిరాంకం}{!!{constant}}" గురించి మాట్లాడవచ్చు. అల్గోరిథం ఉత్పత్తి చేసే వృక్షంలోని ప్రతి సూచికల సీక్వెంట్‌కు \tetoken{స్థిరాంకాల}{!!{constant}s} సమితి అనుబంధంగా ఉంటుంది. ~$0$ దశ సీక్వెంట్‌కు కేటాయించే \tetoken{స్థిరాంకాల}{!!{constant}s} సమితి, దానిలో కనిపించే అన్ని \tetoken{స్థిరాంకాల}{!!{constant}s} సమితి (అవి ఉంటే); \tetoken{స్థిరాంకాలు}{!!{constant}s} లేకపోతే $\{\Obj c_0\}$.

$n$ దశలో అల్గోరిథం సూచికల సీక్వెంట్‌ల వృక్షాన్ని, వాటికి అనుబంధ \tetoken{స్థిరాంకాల}{!!{constant}s} సమితులను ఉత్పత్తి చేసిందనుకుందాం. $n+1$ దశలో ప్రతి అత్యుపరి సీక్వెంట్‌కు కిందివిధంగా చేస్తాం.

సీక్వెంట్ ఎడమ, కుడి రెండు వైపులా \tetoken{ఒక సూత్రం}{!!a{formula}}~$!A$ ఉంటే, లేదా ఎడమవైపున~$\lfalse$ ఉంటే, ఏమీ చేయం: అటువంటి సీక్వెంట్‌లకు~\Log{G3c}లో \tetoken{నిరూపణలు}{!!{proof}s} ఉన్నాయి; ఆ అత్యుపరి సీక్వెంట్‌కు ఇక చేయవలసింది లేదు.

సీక్వెంట్‌లో అణుకాని \tetoken{సూత్రం}{!!{formula}} ఏదీ లేకపోతే అల్గోరిథాన్ని విఫలమై ఆపాలి. $\Gamma \Sequent \Delta$కు నిరూపణ లేదు.

లేకపోతే అతి చిన్న సూచిక~$i$ ఉన్న అణుకాని \tetoken{సూత్రం}{!!{formula}}~$!A$ను ఎంచుకోండి. ఆ అత్యుపరి సీక్వెంట్ $!A^i, \Pi
\Sequent \Lambda$ లేదా $\Pi \Sequent \Lambda, !A^i$ రూపంలో ఉంటుంది. సీక్వెంట్ గరిష్ఠ సూచికకు ఒకటి కలిపితే వచ్చే సంఖ్యను $k$గా, దానికి కేటాయించిన \tetoken{స్థిరాంకాల}{!!{constant}s} సమితిని $C$గా తీసుకుందాం.
\sourcecorrection{OLTEPTPSSAL-001}{మూలం $k$ను గరిష్ఠ సూచికగానే నిర్వచించి కొత్త సూత్రాలకు అదే సూచికను వాడింది; పాత సూచికతో పునరుక్తి రాకుండా $k$ను గరిష్ఠ సూచికకంటే ఒకటి ఎక్కువగా తీసుకున్నాం.}

$!A^i$ను "తగ్గిస్తాం"; అంటే~$!A$ యొక్క \tetoken{ప్రధాన సంయోజకం}{!!{main operator}}, $!A^i$ ఎడమవైపునా కుడివైపునా ఉందనే విషయం నిర్ణయించే నిగమన నియమం ప్రకారం కొత్త అత్యుపరి సీక్వెంట్‌లను ఉత్పత్తి చేస్తాం.

\begin{enumerate}
  \item $!A \ident !B \land !C$ ఎడమవైపున ఉంటే: $(!B \land !C)^i, \Pi \Sequent \Lambda$ పైన ఒక కొత్త అత్యుపరి సీక్వెంట్‌ను చేర్చండి:
  \[
  \Axiom$!B^k, !C^{k+1}, \Pi \fCenter \Lambda$
  \RightLabel{\LeftR{\land}}
  \UnaryInf$(!B \land !C)^i, \Pi \fCenter \Lambda$
  \DisplayProof
  \]
  కొత్త సీక్వెంట్‌కు కేటాయించే \tetoken{స్థిరాంకాల}{!!{constant}s} సమితి~$C$.
  \item $!A \ident !B \land !C$ కుడివైపున ఉంటే: $\Pi \Sequent \Lambda, (!B \land !C)^i$ పైన రెండు కొత్త అత్యుపరి సీక్వెంట్‌లను చేర్చండి:
  \[
  \Axiom$\Pi \fCenter \Lambda, !B^k$
  \Axiom$\Pi \fCenter \Lambda, !C^k$
  \RightLabel{\RightR{\land}}
  \BinaryInf$\Pi \fCenter \Lambda, (!B \land !C)^i$
  \DisplayProof
  \]
  కొత్త సీక్వెంట్‌కు కేటాయించే \tetoken{స్థిరాంకాల}{!!{constant}s} సమితి~$C$.

  \item $!A \ident \lnot !B$, $!A \ident !B \lor !C$, $!A \ident !B \lif !C$ సందర్భాలు ఇలాంటివే; వాటిని అభ్యాసాలుగా వదిలాం.

  \item $!A \ident \lexists[x][!B(x)]$ ఎడమవైపున ఉంటే: $\lexists[x][!B(x)]^i, \Pi \Sequent \Lambda$ పైన ఒక కొత్త అత్యుపరి సీక్వెంట్‌ను చేర్చండి:
  \[
  \Axiom$!B(c)^k, \Pi \fCenter \Lambda$
  \RightLabel{\LeftR{\lexists}}
  \UnaryInf$\lexists[x][!B(x)]^i, \Pi \fCenter \Lambda$
  \DisplayProof
  \]
  ఇక్కడ $c$ అనేది~$C$లో లేని మొదటి \tetoken{స్థిరాంకం}{!!{constant}}. కొత్త సీక్వెంట్‌కు కేటాయించే \tetoken{స్థిరాంకాల}{!!{constant}s} సమితి $C \cup \{c\}$.
  \item $!A \ident \lexists[x][!B(x)]$ కుడివైపున ఉంటే: $\Pi \Sequent \Lambda, \lexists[x][!B(x)]^i$ పైన ఒక కొత్త అత్యుపరి సీక్వెంట్‌ను చేర్చండి:
  \[
  \Axiom$\Pi \fCenter \Lambda, \lexists[x][!B(x)]^k, !B(t)^{k+1}$
  \RightLabel{\RightR{\lexists}}
  \UnaryInf$\Pi \fCenter \Lambda, \lexists[x][!B(x)]^i$
  \DisplayProof
  \]
  ఇక్కడ $t$ అనేది~$C$లోని \tetoken{స్థిరాంకాలను}{!!{constant}s} మాత్రమే కలిగి, ఈ సీక్వెంట్ కింద కుడివైపున $\lexists[x][!B(x)]$ను తగ్గించడానికి ఇంతకు ముందు వాడని మొదటి పదం (అలాంటి పదాలన్నింటినీ ఇప్పటికే వాడితే, ఆ సమితిలోని మొదటి స్థిరాంకాన్ని మళ్లీ వాడండి). కొత్త సీక్వెంట్‌కు కేటాయించే \tetoken{స్థిరాంకాల}{!!{constant}s} సమితి $C$.
  \sourcecorrection{OLTEPTPSSAL-002}{కుడి అస్తిత్వ నియమ వృక్షానికి మూలం \LeftR{\lexists} అనే పేరు పెట్టింది; అది \RightR{\lexists}.}
  \sourcecorrection{OLTEPTPSSAL-003}{కుడి అస్తిత్వ కేసులో తగ్గించే సీక్వెంట్ గద్యానికి మూలం సూచికను వదిలింది; వృక్షం ప్రకారం $i$ సూచికను చేర్చాం.}
  \sourcecorrection{OLTEPTPSSAL-004}{మూలంలోని పునర్వినియోగ ప్రత్యామ్నాయం $\Obj c_0$; కానీ అది $C$లో ఉండాల్సిన అవసరం లేదు. సీక్వెంట్‌లోని అన్ని స్థిరాంకాలూ కేటాయించిన సమితిలో ఉండాలనే తరువాతి వాదనను నిలపడానికి, ఖాళీ కాని $C$లోని మొదటి స్థిరాంకాన్నే మళ్లీ వాడాం.}
  \item $!A \ident \lforall[x][!B(x)]$ సందర్భాలు ఇలాంటివే; వాటిని అభ్యాసాలుగా వదిలాం.
\end{enumerate}

$n+1$ దశలో ఏ అత్యుపరి సీక్వెంట్‌కూ కొత్త సీక్వెంట్‌ను చేర్చకపోతే, అల్గోరిథం విజయవంతంగా ముగుస్తుంది. అప్పుడు $n+1$ దశ వృక్షంలోని సూచికలన్నింటినీ తొలగిస్తే~$\Gamma \Sequent \Delta$కు \tetoken{ఒక నిరూపణ}{!!a{proof}} వస్తుంది. అత్యుపరి సీక్వెంట్‌లన్నీ $!A, \Pi \Sequent \Lambda, !A$ లేదా $\lfalse, \Pi
\Sequent \Lambda$ రూపంలో ఉంటాయి. అనుమితులన్నీ~\Log{G3c}లోని సరైన అనుమితులే. ప్రతిపాదనా అనుమితులకు ఇది స్పష్టం. ప్రతి దశలో సీక్వెంట్‌కు కేటాయించిన \tetoken{స్థిరాంకాల}{!!{constant}s} సమితి~$C$, ఆ సీక్వెంట్‌లో కనిపించే అన్ని \tetoken{స్థిరాంకాలను}{!!{constant}s} కలిగి ఉంటుందని గమనించండి. అందువల్ల \LeftR{\lexists}, \RightR{\lforall} వాడే తగ్గింపుల్లో స్వతంత్ర చరరాశి నియమం తీరుతుంది: స్వతంత్ర చరరాశి~$c$, నిష్కర్షకు కేటాయించిన~$C$లో లేని \tetoken{ఒక స్థిరాంకం}{!!a{constant}} కాబట్టి, $c$ నిష్కర్షలోనూ కనిపించదు. కాబట్టి:

\begin{prop}
\Log{G3c} నిరూపణ అన్వేషణ అల్గోరిథం సరైనది: అది విజయవంతంగా ముగిస్తే, ప్రారంభ సీక్వెంట్~$\Gamma \Sequent \Delta$కు \tetoken{ఒక నిరూపణ}{!!a{proof}} ఉంది.
\end{prop}

\end{document}
